% Epilogue of Verifiable Intelligence.
\specialchapter{Epilogue}{The oversight ledger}{P,V,P,V,P,V,D}\label{ch:epilogue}
\begin{refsection}[chapters/16-epilogue.bib]
\vifloatprefix{E}

Babai's Arthur could toss coins and do a little arithmetic, and the theory that grew from his story asked what such a king can learn from a wizard he cannot trust. This book has asked that question fifteen times, once for each protocol, and has answered it each time in the same three parts: whether honesty is the wizard's best course, what the king must spend to find out, and how often a lie gets past him. Read one at a time, the answers are theorems about particular protocols. Read side by side, they say something about the oversight problem as a whole. This epilogue sets them side by side.

\addsec{The ledger}

\Cref{tab:epilogue:ledger1,tab:epilogue:ledger2} collect the protocol cards of the fifteen chapters. Each row compresses a chapter's answers into a phrase, so the precise statements, with their hypotheses, are in the chapters; the tables are a map of where the guarantees are strong and where they are thin.

\begin{table}[p]
\begin{fullwidth}
\footnotesize
\renewcommand{\arraystretch}{1.18}
\begin{tabularx}{\linewidth}{@{}>{\raggedright\arraybackslash}p{24mm}>{\raggedright\arraybackslash}X>{\raggedright\arraybackslash}X>{\raggedright\arraybackslash}X@{}}
\toprule
\textsc{\color{prover}protocol} & \textsc{\color{verifier}is honesty an equilibrium?} & \textsc{\color{verifier}what does the verifier spend?} & \textsc{\color{verifier}how often is a lie caught?}\\
\midrule
\multicolumn{4}{@{}l}{\textsc{\color{prover}part i. fixed provers}}\\[0.3ex]
1\enspace certificates & yes: valid certificates always pass, lies never; but finding one can be $\cc{NP}$-hard, and a negative claim may need an exponentially long one & one message and no coins; time polynomial in the certificate's length & always, against every prover, relative to the formal statement \\
2\enspace spot checks (Hadamard PCP) & yes for a fixed string; the honest proof has $2^\ell+2^{\ell^2}$ bits, polynomially many with the PCP theorem & $12$ queries and $O(\ell^2+M)$ random bits per run & at least $1/4$ per run, $1-(3/4)^t$ after $t$; never against a prover that sees the queries \\
3\enspace program checking & the correct program pays nothing extra, but a consistent liar does as well as the truth & $O(t)$ calls to the program, plus one per labeled example in self-correction; own work little-oh of the program's & at least $1-2^{-t}$ for products and isomorphism, against every fixed program; never against a consistent liar, nor against memory in self-correction \\[0.4ex]
\multicolumn{4}{@{}l}{\textsc{\color{prover}part ii. conversation}}\\[0.3ex]
4\enspace private coins (graph non-isomorphism) & dominant: truth accepted always, every lie with probability at most $2^{-t}$; the honest prover must decide graph isomorphism & one bit and one permutation per round & at least $1-2^{-t}$ against every prover blind to the coins; never against one that sees them \\
5\enspace sumcheck & a best response, but the honest messages are sums over exponentially many points & $\ell$ rounds, $O(\ell+\sum_jd_j)$ field operations, one evaluation of $g$ & at least $1-\sum_jd_j/\abs{\F}$ against every prover; the best liar survives with probability $1-\prod_j(1-d_j/\abs{\F})$ \\
6\enspace delegation (GKR) & a best response the prover can afford: $O(S\log S)$ operations & $O(D\log S)$ rounds; time linear in the input and output plus $O(D\log S)$, and the wiring evaluations & at least $1-5(D+1)\ceil{\log_2S}/\abs{\F}$ against every prover \\[0.4ex]
\multicolumn{4}{@{}l}{\textsc{\color{prover}part iii. competition}}\\[0.3ex]
7\enspace debate & a minimax strategy and the only equilibrium answer; it can be $\cc{PSPACE}$-hard to play & polynomial time and no coins; in bisection, $\ceil{\log_2T}$ rounds and one step checked & every lie loses under optimal play; no guarantee against a liar who out-computes the honest side \\
8\enspace cross-examination & when $c=1$, a Nash equilibrium and an optimal pair for isolated provers, who must coordinate on one witness & one run of an oracle verifier and one extra question & at least $(1-s)/q$ per round against isolated classical provers; nothing in the clause--variable game if one bit per round passes between them \\
9\enspace efficient debate & affordable ($2T-1$ simulated steps), every equilibrium truthful, for a fixed procedure; maybe not when the debaters choose the decomposition & $\log_2T$ rounds, one question to a human; $O(K^2)$ samples when judgments are noisy & every lie, with deterministic judgments; survival at most $2/5$ per run with noisy ones; at most $(a+1)/N$ detection against obfuscation \\
\bottomrule
\end{tabularx}
\end{fullwidth}
\caption{The ledger, Parts I--III. Each entry summarizes the protocol card of its chapter; the chapters state the hypotheses. Here $t$ is the number of repetitions, $\ell$ the number of variables (of unknowns in \cref{ch:spotchecks}, whose system has $M$ equations), $d_j$ their degrees, $S$ and $D$ a circuit's size and depth, $T$ a computation's length, $c$, $s$ and $q$ an oracle verifier's completeness, soundness error and number of queries, $K$ a Lipschitz constant, and $a$ and $N$ the numbers of checked leaves and of all leaves of an obfuscated argument (\cref{ch:efficientdebate}).}\label{tab:epilogue:ledger1}
\end{table}

\begin{table}[p]
\begin{fullwidth}
\footnotesize
\renewcommand{\arraystretch}{1.18}
\begin{tabularx}{\linewidth}{@{}>{\raggedright\arraybackslash}p{24mm}>{\raggedright\arraybackslash}X>{\raggedright\arraybackslash}X>{\raggedright\arraybackslash}X@{}}
\toprule
\textsc{\color{prover}protocol} & \textsc{\color{verifier}is honesty an equilibrium?} & \textsc{\color{verifier}what does the verifier spend?} & \textsc{\color{verifier}how often is a lie caught?}\\
\midrule
\multicolumn{4}{@{}l}{\textsc{\color{prover}part iv. incentives}}\\[0.3ex]
10\enspace rational proofs & the true report is strictly dominant for a risk-neutral prover, but by a margin of only $2\cdot4^{-\ell}$; once computation costs anything, a cheap guess pays better unless the payment grows like $4^\ell$ & one message, $\ell$ random bits, one evaluation, a budget of $2$ & never: a lie of size $\Delta$ costs $2(\Delta/2^\ell)^2$ in expectation, and spotting it in the payments takes about $4^\ell/\Delta^2$ runs \\
11\enspace prover--verifier games & only when the verifier commits before the instance is drawn are the equilibria exactly the complete and sound pairs; other orders have bad equilibria; in practice honesty costs accuracy & one evaluation of a learned verifier; the real cost is training & every lie of a prover in the class, at an exact equilibrium; at an approximate one, less often on instances rare in training; beyond the class, only as measured \\
12\enspace PAC verification & a best response, not strict: small lies pass but move the output little & $\ceil{4\ln(2/\delta)/\eps}$ trusted labels, cheap unlabeled points, one fit & the output is bad with probability at most $\delta$; checks the prover can recognize never catch anything \\[0.4ex]
\multicolumn{4}{@{}l}{\textsc{\color{prover}part v. secrets and limits}}\\[0.3ex]
13\enspace zero knowledge & a best response that discloses nothing; the prover needs the certificate & polynomial time; about $0.69k\abs{E}$ rounds for error $2^{-k}$ & at least $1/\abs{E}$ per round ($1/2$ for graph isomorphism), against every prover if the commitments are perfectly binding \\
14\enspace proofs of training & honest training passes but is not enforced: a spoof passes when the learning rate is free; refereed training makes honesty a best response to an honest rival & a fraction $q/N$ of the run replayed, $O(q\log N)$ hashes; a referee needs $\ceil{\log_2T}$ rounds and one step & at least $1-(1-\eps)^q$, less a negligible term, against bounded developers; about $q/N$ for a wrong model; never against logs valid within the tolerance; refereed, always, unless a hash collision is found \\
15\enspace backdoor audits & not an equilibrium: a backdoored model passes as often as an honest one, up to a negligible amount; mitigation makes the backdoor worthless & no polynomial number of queries helps the audit; mitigation costs $2t$ queries per input & detection at most the false-alarm rate plus a forger's advantage; after mitigation a trigger works with probability at most $e^{-2t(1/2-2\eps)^2}$ \\
\bottomrule
\end{tabularx}
\end{fullwidth}
\caption{The ledger, Parts IV--V. Here $\ell$ is the number of variables of the counted formula, $\Delta$ a lie's size, $\delta$ and $\eps$ the confidence and accuracy of a learning guarantee ($\eps$ is the fraction of invalid segments of a log in \cref{ch:training}, and the fraction of inputs on which the model is wrong in \cref{ch:backdoors}), $\abs{E}$ the number of edges of a graph, $q$ and $N$ the numbers of checked and logged segments, $T$ the number of training steps, and $t$ the number of self-corrected trials.}\label{tab:epilogue:ledger2}
\end{table}

\addsec{What the ledger shows}

\paragraph{The price of honesty.} The left column of the ledger is where the oversight problem differs most from the problem that complexity theory solved. Every protocol on the upper rungs of the prologue's ladder makes honesty a best response only for a prover that can compute the honest strategy, and the honest strategy is as hard as the claim: sums over exponentially many points (\cref{ch:sumcheck}), perfect play of a game of polynomial length (\cref{ch:debate}), a witness of exponential size shared by two provers (\cref{ch:crossexam}). A theorem that the truth would win against an optimal liar says nothing about a model that can afford only part of the truth. The protocols that protect a bounded model are those whose honest strategy costs about what the task costs: delegation of a computation the model has already run (\cref{ch:delegation}), debate over an agreed procedure whose steps the model has already taken (\cref{ch:efficientdebate}), labeling points the model can label (\cref{ch:pacverify}), proving a claim whose certificate the model holds (\cref{ch:certificates,ch:zeroknowledge}). The arc from \cref{ch:sumcheck} to \cref{ch:delegation}, and from \cref{ch:debate} to \cref{ch:efficientdebate}, is the most important movement in the book.

\paragraph{The verifier's secret.} The right column fails in the same way in six chapters: a prover that can see or predict the verifier's randomness, or tell its checks apart, is never caught. The adaptive prover defeats spot checks (\cref{ch:spotchecks}), a program with memory defeats self-correction (\cref{ch:consistency}), the prover that sees the coins defeats the non-isomorphism protocol (\cref{ch:interaction}), checks that the prover can recognize catch nothing (\cref{ch:pacverify}), a developer that knew the checks before committing to its log could make every checked segment valid (\cref{ch:training}), and mitigation works only because each of the self-corrector's queries is uniform, whatever the input that matters (\cref{ch:backdoors}). In these protocols the verifier's advantage over a stronger prover is information the prover does not have, and the questions it asks after the prover has committed. In AI oversight the corresponding threats are contamination of evaluations, models that recognize when they are being tested, and protocols in which the model can choose what is checked. Commitment before challenge, which the theory takes for granted, has to be engineered.

\paragraph{Isolation and competition.} Two provers are stronger than one only if they cannot coordinate (\cref{ch:crossexam}), and a debate protects only if the honest debater can find the flaw (\cref{ch:efficientdebate}). Copies of one model share their weights, and debaters trained against each other share a training run, so the conditions under which competition helps are conditions on how the competitors are built, not only on how they are questioned.

\paragraph{Incentives are not detection.} The middle of the book replaced the worst-case adversary with a prover that is paid or trained. The guarantees change kind. A rational proof never catches a lie; it makes the lie cost money in expectation, and the cost can be smaller than the cost of computing the truth (\cref{ch:rational}). A prover--verifier game makes every equilibrium a complete and sound protocol only for one order of moves, and its soundness beyond the provers seen in training is measured rather than proved (\cref{ch:pvg}). These are the kinds of guarantee that training a model can supply; they are weaker than the theorems of Parts I--III, and the ledger shows exactly where.

\paragraph{The limit.} If one-way functions exist, some lies escape every efficient audit that queries the model, whatever questions it asks, and, for some models and under stronger assumptions, every audit that reads the weights (\cref{ch:backdoors}). The response the theory offers is not better detection but a change of the game: a verifier that cannot detect a backdoor can sometimes make it useless, by asking the model only questions whose answers it can correct. Oversight that cannot certify the prover can still certify its own use of the prover.

\addsec{What the theory asks of practice}

The guarantees of this book transfer to oversight under conditions that are easy to state and hard to meet. The claim must be stated formally, and the formal statement must say what the overseer means. The verifier's randomness must be fresh, secret, and drawn after the prover has committed. Copies of a model that are questioned apart must be unable to coordinate. The honest strategy must be one the model can afford. Human judgment should enter as answers to small, well-posed questions, with a model of how those answers err. And the prover's commitments, to a log, to a model's weights, to an answer, must be binding before the checks begin. Each chapter's section on what the model leaves out is a catalog of the ways a real deployment can violate one of these conditions.

\addsec{Open problems}

The chapters end with open problems, and some of them recur. Do all efficiently decidable claims have interactive proofs in which the verifier runs in nearly linear time and the honest prover in polynomial time, or only those decidable in small depth or small space (\cref{ch:delegation})? What is the right model of systematic human error in protocols that consult a human, and which protocols survive it (\cref{ch:efficientdebate})? By a 2026 preprint, stability is necessary to defeat obfuscated arguments when people are only queried as a black box; what other structure of human judgment, such as robustness to a few wrong answers, would do instead (\cref{ch:efficientdebate,ch:delegation})? Which multi-prover protocols keep a useful soundness bound when copies of one model can exchange a few bits, and which channels between copies does a deployment leave open (\cref{ch:crossexam})? Are there error-correcting encodings of informal arguments, which people could spot-check as a verifier spot-checks a PCP (\cref{ch:spotchecks})? When does gradient training find the verifier-leading equilibria that make prover--verifier games sound, beyond the small channel for which \cref{ch:pvg} proves that it does? Which questions about a predicate that the verifier can only sample, such as a human judgment, admit rational proofs with a noticeable reward gap and a polynomial budget (\cref{ch:rational})? Is there a proof of training whose soundness rests on stated assumptions rather than on the failure of known attacks (\cref{ch:training})? Which classes of models can be cleaned of backdoors without detecting them (\cref{ch:backdoors})? And can consistency checks for natural-language tasks be given soundness guarantees (\cref{ch:consistency})?

\addsec{Arthur's tools}

None of the protocols in this book makes Merlin honest. They make honesty his best course, or make a lie expensive, or make it detectable, or make it useless, and each does so only under conditions the chapters state. That is less than one might want from a theory of oversight, and more than one has without it. When people oversee systems more capable than themselves, they will be in Arthur's position whether or not they have read Babai: deciding, with little computation and some coins, whether to believe a stronger party. The theory of interactive proofs is the most developed account of what can be asked of that party, and of what it costs to ask. It does not say whether the systems we build will behave like the provers it assumes. That is a question about the systems, and it is the question that the next textbooks in this subject will have to answer.

\notesandrefs

The ledger summarizes the chapters, and the references are there. Two surveys are worth reading after this book: \textcite{Babai1990}, a participant's account of the discoveries of 1989 and 1990, and \textcite{Goldreich2018} on doubly-efficient interactive proofs, which is the theory's own account of the price of honesty. The research agenda of the UK AI Security Institute's alignment team \parencite{HiltonEtAl2025} lists open problems at the boundary between this book's theory and the practice of AI oversight.

\printchapterbib
\end{refsection}
\vifloatrestore
